% !TeX program = xelatex

\documentclass[14pt]{extarticle}
\usepackage[drawstamps=true]{bstunote}
\usepackage[a4paper, left=23mm, right=9mm, top=9mm, bottom=24mm]{geometry}
\usepackage[english, russian]{babel}

% Включаем гиперссылки в оглавлении
\usepackage[colorlinks=true,linkcolor=blue,urlcolor=blue,citecolor=green!70!black,hypertexnames=false]{hyperref}

\usepackage{lipsum}

\StampAuthor{Иванов И.И.}{3.06.26\,}
\StampCheckBy{Осипов О.В.}{3.06.26\,}
\StampNormControl{Бондаренко Т.В.}{11.06.26\,}
\StampApprovedBy{Поляков В.М.}{25.06.26\,}
\StampWorkTheme{Разработка программного обеспечения для прогнозирования погоды на основе анализа спутниковых снимков с использованием методов машинного обучения}
\StampOrganization{БГТУ~им.~В.Г.~Шухова, \makebox{ПВ-222}}

% Шрифт в рамках
\usepackage{fontspec}
\newfontfamily\GOSTFontFamily[
FakeSlant=0.3 	% Используем искусственный наклон
]{GOST type A}  % Название шрифта, который будет в рамках
\StampFont{\GOSTFontFamily}

\setmainfont{PT Astra Serif} 	 % Шрифт с засечками
\setsansfont{PT Astra Sans}		 % Шрифт без засечек
\setmonofont[Scale=0.9]{PT Mono} % Моноширинный шрифт

\begin{document}

\initpagenumbering

\section*{OПРЕДЕЛЕНИЯ, СOКРAЩЕНИЯ И OБOЗНAЧЕНИЯ}
\par РФ -- Российская Федерация.
\par СВО -- специальная военная операция.
\par ВУЗ -- высшее учебное заведение.
\par API (Application Programming Interface) -- программный интерфейс, обеспечивающий взаимодействие между компонентами программного обеспечения.
\par RGB (red, green, blue) -- цветовой формат машинной графики, основанный на смешении компонентов красного, зелёного и синего цветов.
\par WWW (World Wide Web) -- всемирная паутина, система взаимосвязанных гипертекстовых документов.

% Содержание
\renewcommand{\contentsname}{\rucontentsname}
\tableofcontents

\section*[ВВЕДЕНИЕ]{ВВЕДЕНИЕ}
\lipsum[1-10]
\section{АНАЛИЗ СОСТОЯНИЯ ИССЛЕДОВАНИЙ И РАЗРАБОТОК В ОБЛАСТИ ПРОГНОЗИРОВАНИЯ ПОГОДЫ}
\lipsum[11-20]
\section{РАЗРАБОТКА АРХИТЕКТУРЫ ПРОГРАММНОГО ОБЕСПЕЧЕНИЯ ДЛЯ ПРОГНОЗИРОВАНИЯ ПОГОДЫ}
\lipsum[21-30]
\section{РАЗРАБОТКА АЛГОРИТМОВ И ПРОГРАММНОГО ОБЕСПЕЧЕНИЯ ДЛЯ ПРОГНОЗИРОВАНИЯ ПОГОДЫ}
\lipsum[31-40]
\section*[ЗАКЛЮЧЕНИЕ]{ЗАКЛЮЧЕНИЕ}
\lipsum[41-50]

\newpage
% Оформление списка литературы
\renewcommand{\refname}{\rubibliographyname} % Заголовок
\phantomsection
\addcontentsline{toc}{section}{\rubibliographyname}% Добавим список литературы в оглавление
\begin{thebibliography}{99}
	\bibitem{gost19201}
	\textbf{ГОСТ 19.201-78.} Техническое задание. Требования к содержанию и оформлению.  -- Введ. 1980-01-01. -- Москва : Стандартинформ, 2010. --~4~с.
\end{thebibliography}

% ПРИЛОЖЕНИЯ
\appendix

\section{Блок-схемы основных алгоритмов прогнозирования погоды}
\lipsum[1-10]
\section{Листинги разработанных подпрограмм для прогнозирования погоды}
\lipsum[11-20]
\end{document}
